\chapter{Reproducible environments with pixi}
\label{ch:pixi}

\chaptergoals{4 of 6 in Part~I.}
{\cref{ch:git}: you can commit. A network connection to download packages.}
{create a project environment from a manifest and a lock file, run commands
and tasks in it, add packages and a second environment, and recognise how the
ShipSoft repositories use all of this.}

\section{Why an environment, and why pixi}

``It works on my machine'' is the oldest problem in shared software. An
analysis depends on Python, ROOT, and everything ROOT depends on; the SHiP
packages add Geant4, GeoModel, Phlex and compilers. Two people with different
versions of any of these can get different results, or none.

\href{https://pixi.sh}{pixi} solves this per project. It installs packages
from conda channels (\texttt{conda-forge} for the general HEP stack,
\href{https://prefix.dev/channels/ship}{\texttt{prefix.dev/ship}} for the SHiP
packages) into a directory inside the project, \file{.pixi/}. Nothing is
installed system-wide, nothing needs administrator rights, and every project
has exactly the versions it asks for. Every ShipSoft repository is a pixi
project, and so is the team repository of \cref{ch:team}.

A pixi project is described by two files, and both are committed to git:

\begin{center}\small
\begin{tabularx}{\linewidth}{l L L}
\toprule
 & \file{pixi.toml}, the \emph{manifest} & \file{pixi.lock}, the \emph{lock file} \\
\midrule
written by & you & pixi, never by hand \\
says       & what the project needs: \texttt{root\_base >=6.36} & exactly what was installed:
                                                             every package, version and build, per platform \\
changes    & when you change the requirements & when pixi solves again (\cmd{pixi add}, \cmd{pixi update}) \\
gives      & the recipe & the same environment for everybody, and for CI \\
\bottomrule
\end{tabularx}
\end{center}

\section{Installing pixi}

\begin{shell}
curl -fsSL https://pixi.sh/install.sh | bash
# open a new terminal, then:
pixi --version
\end{shell}

\begin{pitfall}
pixi keeps downloaded packages in a cache (\cmd{pixi info} shows where) and
the environments in each project's \file{.pixi/}. Both can take several
gigabytes. On machines with a small home quota, such as lxplus, point the
cache at a larger disk with \cmd{export PIXI\_CACHE\_DIR=<somewhere big>} in
your shell's start-up file. \tbd{recommended location on the tutorial machines}
\end{pitfall}

\section{Your first workspace}

\begin{exercise}{An environment with PyROOT}
\label{ex:pixifirst}
Next to your git playground:
\begin{shell}
mkdir pixi-playground && cd pixi-playground
git init
pixi init                                  # writes pixi.toml, .gitignore, .gitattributes
pixi add python "root_base>=6.36"          # solves, writes pixi.lock, installs
pixi run python -c "import ROOT; print(ROOT.gROOT.GetVersion())"
\end{shell}
Open \file{pixi.toml}: what did \cmd{pixi add} write into it? Look at the size
of \file{pixi.lock} and search it for \texttt{root\_base}: which exact version
and build did you get? What does the new \file{.gitignore} keep out of git?
Commit the project:
\begin{shell}
git add pixi.toml pixi.lock .gitignore .gitattributes
git commit -m "feat: environment with PyROOT"
\end{shell}
\end{exercise}

\begin{checkpoint}
A ROOT version of 6.36 or later printed, and one commit containing the four
files, but not \file{.pixi/}.
\end{checkpoint}

There are two ways to use the environment:

\begin{description}[style=nextline]
  \item[\cmd{pixi run <command>}] runs one command inside the environment
    and returns. Use it in scripts, in CI, and whenever you can.
  \item[\cmd{pixi shell}] opens a shell with the environment active until you
    \cmd{exit}. Convenient for interactive work; leave it before you enter
    the shell of another project, or you get a mixture of both.
\end{description}

\section{Tasks}

Typing the same long commands is error-prone, and the next person will type
them differently. A \emph{task} gives a command a name, stored in the
manifest, so that everybody, including CI, runs exactly the same thing.

\begin{exercise}{Name your commands}
\label{ex:pixitasks}
\begin{shell}
echo 'import ROOT; print("hello from ROOT", ROOT.gROOT.GetVersion())' > hello.py
pixi task add hello "python hello.py"
pixi run hello
pixi task list
\end{shell}
Look at the \texttt{[tasks]} table in \file{pixi.toml}. Tasks can depend on
other tasks and take arguments; edit the table by hand to:
\begin{tomlcode}
[tasks]
hello = "python hello.py"
greet = { cmd = "echo Hello, {{ name }}", args = ["name"] }
all = { cmd = "echo done", depends-on = ["hello"] }
\end{tomlcode}
and try \cmd{pixi run greet SHiP} and \cmd{pixi run all}. Commit.
\end{exercise}

\section{More than one environment}

Some tools are needed only for some jobs: a linter, a benchmark tool, a
documentation generator. pixi groups them into \emph{features} and combines
features into named \emph{environments}, so that the linter does not end up
in the analysis environment.

\begin{exercise}{A separate environment for the linter}
\begin{shell}
pixi add --feature lint ruff
pixi workspace environment add lint --feature lint --no-default-feature
pixi task add --feature lint lint "ruff check ."
pixi run lint                  # pixi picks the environment that has the task
pixi run -e lint ruff --version
\end{shell}
Read the \texttt{[feature.lint]} and \texttt{[environments]} tables that
appeared in \file{pixi.toml}. Write an unused \texttt{import os} into
\file{hello.py} and run \cmd{pixi run lint} again. Remove it, and commit.
\end{exercise}

\section{Keeping everybody on the same versions}

\begin{description}[style=nextline]
  \item[\cmd{pixi install -{}-locked}] installs exactly what \file{pixi.lock}
    says, and fails if the lock no longer matches the manifest. This is what
    CI runs: a change of requirements must come with its lock file.
  \item[\cmd{pixi update}] solves again and moves to newer versions within
    the manifest's ranges. The new \file{pixi.lock} is a change like any
    other: commit it, on its own, so that its effect can be tested and
    reviewed. In the ShipSoft repositories a bot opens such a pull request
    every month (\cref{ch:ci}).
  \item[\cmd{pixi list}] shows the installed packages and versions.
\end{description}

\begin{pitfall}
Never edit \file{pixi.lock} by hand, and do not resolve conflicts inside it
(the \file{.gitattributes} that \cmd{pixi init} writes tells git not to
try). If a rebase conflicts on \file{pixi.lock}, take the version of
\texttt{main} (\cmd{git checkout upstream/main -{}- pixi.lock}), run
\cmd{pixi install} to solve again with your manifest, and \cmd{git add} the
result. Better still: do not change the environment in a pull request that is
about something else.
\end{pitfall}

\section{How the ShipSoft repositories use pixi}

Every repository has a \file{pixi.toml} and a \file{pixi.lock} at its top.
Their channels are the SHiP channel on \texttt{prefix.dev} and
\texttt{conda-forge}.
Building and testing are tasks, chained with \texttt{depends-on}, so that one
command does everything that is needed. From \repo{Shannon}:
\begin{tomlcode}
build = { cmd = "cmake --build build -j8", depends-on = ["configure"] }
install = { cmd = "cmake --install build", depends-on = ["build"] }
test = { cmd = "ctest --test-dir build --output-on-failure", depends-on = ["build"] }
smoke = { cmd = "scripts/smoke.sh", depends-on = ["build"] }
\end{tomlcode}
so \cmd{pixi run test} configures, builds and tests a fresh checkout.

\begin{itemize}
  \item \textbf{A \texttt{lint} environment} holds the formatting and lint
    tools (\texttt{prek}, \texttt{ruff}, \texttt{clang-format},
    \texttt{codespell}, \texttt{commitizen}, \ldots).
    \cmd{pixi run lint} checks every file; \cmd{pixi run install-hooks}
    installs git hooks that run the same checks at every commit.
  \item \textbf{Other environments} for special jobs, such as \texttt{bench}
    for benchmarks in \pkg{aegir}.
  \item \textbf{A commented-out \texttt{dev} feature} builds a dependency
    from a local checkout, for changing two packages at once (see the
    developer guide of the SHiP documentation). It stays commented out in
    commits, together with the lock-file changes it causes.
  \item \textbf{CI runs the same tasks}, with \cmd{-{}-locked}
    (\cref{ch:ci}). If \cmd{pixi run -{}-locked test} passes on your machine,
    the build check passes in CI.
\end{itemize}

\begin{levelcheck}
\begin{itemize}
  \can say what belongs in \file{pixi.toml} and what in \file{pixi.lock}, and
       why both are committed;
  \can create a workspace, add packages, and run a command in it with
       \cmd{pixi run} or \cmd{pixi shell};
  \can define tasks, and a second environment from a feature;
  \can reproduce an environment exactly with \cmd{-{}-locked}, and update it
       on purpose.
\end{itemize}
\end{levelcheck}

\begin{solution}
The commands of this chapter were checked with pixi~0.82: \cmd{init},
\cmd{task add}, task arguments and \texttt{depends-on}, \cmd{add -{}-feature},
\cmd{workspace environment add -{}-feature -{}-no-default-feature}, automatic
environment selection for a feature-only task, and \texttt{PIXI\_CACHE\_DIR}.
Solving \texttt{python} and \texttt{root\_base} needs network access to
\texttt{conda-forge} and was not run in the review environment: do
\cref{ex:pixifirst} once on the tutorial machines.
\end{solution}
